% GHI to GII (Global Horizontal to Global Inclined Irradiance) Transposition
% Implementation documentation for ghi_to_gii.py

\documentclass[11pt,a4paper]{article}

\usepackage{amsmath, amssymb}
\usepackage{siunitx}
\usepackage{geometry}
\usepackage{graphicx}
\usepackage{tikz}
\usepackage{hyperref}
\usetikzlibrary{shapes.geometric, arrows.meta, positioning}

\geometry{margin=2.5cm}
\sisetup{detect-all}

\title{GHI to GII Transposition Model\\(Implementation Documentation)}
\author{Internal Documentation}
\date{\today}

\begin{document}
\maketitle

\section{Overview}

The function \texttt{ghi\_to\_gii} implements a physically-based transposition model to convert
\emph{global horizontal irradiance} (GHI, \(G_h\)) into \emph{global irradiance on an inclined plane} (GII, \(G_T\)).

Main components:
\begin{itemize}
  \item Solar geometry (declination, hour angle, zenith angle).
  \item Extraterrestrial irradiance and clearness index \(k_t\).
  \item Erbs diffuse fraction correlation.
  \item Hay--Davies diffuse transposition model on the tilted plane.
  \item Direct (beam) component on the tilt.
  \item Ground-reflected component.
  \item Time zone and longitude corrections based on asset configuration.
  \item Numerical stabilisation at very low solar elevation.
\end{itemize}

Inputs (per timestamp):
\begin{itemize}
  \item \(G_h\): measured global horizontal irradiance \([\si{W/m^2}]\).
  \item \(t\): timestamp as timezone-aware datetime.
  \item \(\varphi\): site latitude in degrees (north positive).
  \item \(\lambda\): site longitude in degrees (east positive).
  \item \(\beta\): tilt angle in degrees (0 = horizontal).
  \item \(\gamma\): surface azimuth in degrees (0 = south, east negative, west positive).
  \item \(h\): site altitude [m]. Taken from \texttt{asset\_list.altitude\_m}; default 0.
  \item \(\rho\): ground albedo (0--1). Taken from \texttt{asset\_list.albedo}; default 0.2.
  \item \(T_{\text{local}}\): time zone of the asset (e.g. +08:00) as a fixed offset.
\end{itemize}

Output:
\begin{itemize}
  \item \(G_T\): global irradiance on the tilted plane \([\si{W/m^2}]\).
\end{itemize}

\section{Time Zones and Solar Time}

\subsection{Local Time and Time Zone Offset}

Given an input timestamp \(t\) (typically in UTC or with any timezone), we first convert it to the asset's local timezone \(T_{\text{local}}\):
\[
  t_{\text{loc}} = t \text{ converted to } T_{\text{local}}.
\]

Let the UTC offset of \(T_{\text{local}}\) be
\[
  \Delta_{\text{tz}} = \frac{\text{offset in seconds}}{3600} \quad [\text{hours}].
\]

We define the local civil time in hours:
\[
  h_{\text{loc}} = H + \frac{M}{60} + \frac{S}{3600},
\]
where \(H, M, S\) are the hour, minute, and second of \(t_{\text{loc}}\).

\subsection{Approximate Solar Time and Hour Angle}

Longitude \(\lambda\) is in degrees, positive east.
We approximate solar time as:
\[
  h_{\text{sol}} \approx h_{\text{loc}} + \left(\frac{\lambda}{15} - \Delta_{\text{tz}}\right),
\]
where:
\begin{itemize}
  \item \(\lambda/15\) is the ``longitude time'' in hours,
  \item \(\Delta_{\text{tz}}\) is the timezone offset in hours.
\end{itemize}

Hour angle \(\omega\) (degrees) is then:
\[
  \omega^\circ = 15\,(h_{\text{sol}} - 12),
  \qquad
  \omega = \omega^\circ \cdot \frac{\pi}{180}.
\]

\section{Solar Geometry}

\subsection{Day of Year and Declination}

Let \(N\) be the day of year (1--366), computed from \(t_{\text{loc}}\) including leap years.

Solar declination \(\delta\) (radians) is approximated by:
\[
  \delta = 23.45^\circ \sin\left( \frac{360^\circ}{365} (284 + N) \right),
\]
converted to radians in code:
\[
  \delta_{\text{rad}} = 23.45 \cdot \frac{\pi}{180}
  \cdot
  \sin\left( \frac{360}{365} (284 + N) \cdot \frac{\pi}{180} \right).
\]

Latitude is converted to radians:
\[
  \varphi_{\text{rad}} = \varphi \cdot \frac{\pi}{180}.
\]

\subsection{Solar Zenith Angle}

We compute:
\[
  \cos\theta_z = \sin\varphi_{\text{rad}} \sin\delta_{\text{rad}}
  + \cos\varphi_{\text{rad}} \cos\delta_{\text{rad}} \cos\omega.
\]

If \(\cos\theta_z \le 0\), the sun is below the horizon and
\[
  G_T = 0.
\]

Otherwise, clamp:
\[
  \cos\theta_z \leftarrow \min(1, \cos\theta_z).
\]

Solar elevation angle:
\[
  \theta_e = 90^\circ - \theta_z = 90^\circ - \arccos(\cos\theta_z) \cdot \frac{180}{\pi}.
\]

To avoid numerical artefacts at very low sun angles, we use a minimum elevation threshold \(\theta_{\min}\) (e.g. \(\theta_{\min} = 3^\circ\)):
\[
  \text{if } \theta_e < \theta_{\min} \text{ then } G_T = G_h \quad\text{(pass through measured horizontal)}.
\]

\section{Extraterrestrial Horizontal Irradiance and Clearness Index}

Solar constant:
\[
  G_{sc} = 1367\ \si{W/m^2}.
\]

Eccentricity correction factor:
\[
  E_0 = 1 + 0.033 \cos\left( \frac{360 N}{365} \cdot \frac{\pi}{180} \right).
\]

Extraterrestrial irradiance on horizontal:
\[
  G_{0h} = G_{sc} E_0 \cos\theta_z.
\]
If \(G_{0h} \le 0\), we set \(G_T = 0\).

Clearness index:
\[
  k_t = \frac{G_h}{G_{0h}},
\]
clamped to:
\[
  k_t \leftarrow \max(0, \min(1.2, k_t)).
\]

\section{Erbs Diffuse Fraction}

Diffuse fraction \(k_d\) via the Erbs correlation:

\[
k_d =
\begin{cases}
  1 - 0.09 k_t, & k_t \le 0.22,\\[6pt]
  0.9511 - 0.1604 k_t + 4.388 k_t^2 - 16.638 k_t^3 + 12.336 k_t^4, & 0.22 < k_t \le 0.8,\\[6pt]
  0.165, & k_t > 0.8.
\end{cases}
\]

Diffuse horizontal:
\[
  D_h = k_d G_h.
\]

Beam normal irradiance (direct normal, DNI):
\[
  \text{DNI} = \frac{G_h - D_h}{\cos\theta_z},
\]
with non-negativity:
\[
  \text{DNI} = \max(0, \text{DNI}).
\]

\subsection{Altitude (pressure) correction for beam}

Best practice: at higher altitude, atmospheric path length is shorter, so the direct (beam) irradiance is increased. We use the barometric formula for relative pressure:
\[
  \frac{P}{P_0} = \left(1 - 2.25577\times 10^{-5}\, h\right)^{5.25588},
\]
with \(h\) in metres and \(P_0 = 101325\ \si{Pa}\). Altitude \(h\) is clamped to \([0, 15000]\ \si{m}\). The beam is then scaled:
\[
  \text{DNI} \leftarrow \text{DNI} \cdot \left(\frac{P_0}{P}\right)^{0.7}.
\]
The exponent \(0.7\) is a common broadband approximation. This corrected DNI is used in the Hay--Davies anisotropy and in the beam-on-tilt component.

\section{Incidence Angle on Tilted Plane}

Let:
\[
  \beta = \text{tilt in radians}, \quad
  \gamma = \text{surface azimuth in radians}.
\]

We use the standard incidence-angle formula (orthodox form for surface tilt and azimuth).
In code:

\begin{align*}
  \beta  &= \beta^\circ \cdot \frac{\pi}{180},\\
  \gamma &= \gamma^\circ \cdot \frac{\pi}{180},
\end{align*}

\[
\cos\theta_i =
    \sin\delta_{\text{rad}} \sin\varphi_{\text{rad}} \cos\beta
  - \sin\delta_{\text{rad}} \cos\varphi_{\text{rad}} \sin\beta \cos\gamma
  + \cos\delta_{\text{rad}} \cos\varphi_{\text{rad}} \cos\omega \cos\beta
  + \cos\delta_{\text{rad}} \sin\varphi_{\text{rad}} \cos\omega \sin\beta \cos\gamma
  + \cos\delta_{\text{rad}} \sin\omega \sin\beta \sin\gamma.
\]

Clamp:
\[
  \cos\theta_i \leftarrow \max(0, \cos\theta_i).
\]

\section{Hay--Davies Diffuse on Tilt}

Anisotropy index:
\[
  a = \frac{\text{DNI}}{G_{sc} E_0}, \quad
  a \leftarrow \min(1, \max(0, a)).
\]

Geometric factor:
\[
  R_b = \frac{\cos\theta_i}{\cos\theta_z}.
\]

Diffuse on tilt:
\[
  D_T = D_h \left( a R_b + (1-a)\frac{1 + \cos\beta}{2} \right).
\]

\section{Beam and Ground-Reflected Components}

Beam on tilt:
\[
  B_T = \text{DNI} \cos\theta_i.
\]

Ground-reflected:
\[
  R_T = \rho\,G_h \frac{1 - \cos\beta}{2}.
\]

\section{Final GII}

Total irradiance on the tilt:
\[
  G_T = B_T + D_T + R_T.
\]

We then ensure non-negativity:
\[
  G_T \leftarrow \max(0, G_T).
\]

\section{Special Cases and Filters}

\begin{itemize}
  \item If \(G_h < 0\) or NaN, we return \(0\).
  \item If \(\cos\theta_z \le 0\) (sun below horizon), we return \(0\).
  \item If \(G_{0h} \le 0\), we return \(0\).
  \item If solar elevation \(\theta_e < \theta_{\min}\) (e.g. \(3^\circ\)), we \emph{bypass} the model and set \(G_T = G_h\) to avoid numerical spikes at the horizon with otherwise clean measurements.
\end{itemize}

\section{Flow Chart}

The following simple flowchart illustrates the main control flow per timestamp.

\tikzset{
  block/.style = {rectangle, draw, rounded corners, text width=8em, align=center, minimum height=2em},
  io/.style    = {trapezium, trapezium left angle=70, trapezium right angle=110, draw, align=center},
  decision/.style = {diamond, draw, aspect=1.6, align=center},
  line/.style  = {draw, -{Latex[length=2mm]}}
}

\begin{figure}[p]
\centering
\resizebox{!}{0.8\textheight}{%
\begin{tikzpicture}[node distance=1.8cm]
  \node[io] (input) {Input \\ $G_h, t, \varphi, \lambda, \beta, \gamma, h, \rho, T_{\text{local}}$};
  \node[block, below=of input] (localtime) {Convert $t$ to $t_{\text{loc}}$ using $T_{\text{local}}$};
  \node[block, below=of localtime] (solar) {Compute $N$, $\delta$, $\cos\theta_z$, $\theta_e$};
  \node[decision, below=of solar] (belowHorizon) {$\cos\theta_z \le 0$?};
  \node[block, right=3.8cm of belowHorizon] (return0a) {Set $G_T=0$};
  \node[decision, below=of belowHorizon] (lowElev) {$\theta_e < \theta_{\min}$?};
  \node[block, right=3.8cm of lowElev] (passThrough) {Set $G_T = G_h$};
  \node[block, below=of lowElev] (extrac) {Compute $G_{0h}$, $k_t$};
  \node[decision, below=of extrac] (g0neg) {$G_{0h} \le 0$?};
  \node[block, right=3.8cm of g0neg] (return0b) {Set $G_T=0$};
  \node[block, below=of g0neg] (erbs) {Erbs: $k_d$, $D_h$, DNI};
  \node[block, below=of erbs] (aoi) {Compute $\cos\theta_i$, $R_b$};
  \node[block, below=of aoi] (hay) {Hay--Davies: $D_T$};
  \node[block, below=of hay] (beamground) {Compute $B_T$, $R_T$};
  \node[block, below=of beamground] (sum) {$G_T = B_T + D_T + R_T$};
  \node[block, below=of sum] (clamp) {Clamp $G_T \ge 0$};
  \node[io, below=of clamp] (output) {Output $G_T$};

  \path[line] (input) -- (localtime);
  \path[line] (localtime) -- (solar);
  \path[line] (solar) -- (belowHorizon);
  \path[line] (belowHorizon.east) -- node[above]{Yes} (return0a.west);
  \path[line] (return0a) |- (output);
  \path[line] (belowHorizon) -- node[left]{No} (lowElev);
  \path[line] (lowElev.east) -- node[above]{Yes} (passThrough.west);
  \path[line] (passThrough) |- (output);
  \path[line] (lowElev) -- node[left]{No} (extrac);
  \path[line] (extrac) -- (g0neg);
  \path[line] (g0neg.east) -- node[above]{Yes} (return0b.west);
  \path[line] (return0b) |- (output);
  \path[line] (g0neg) -- node[left]{No} (erbs);
  \path[line] (erbs) -- (aoi);
  \path[line] (aoi) -- (hay);
  \path[line] (hay) -- (beamground);
  \path[line] (beamground) -- (sum);
  \path[line] (sum) -- (clamp);
  \path[line] (clamp) -- (output);
\end{tikzpicture}
}%
\caption{Per-timestamp GHI \(\to\) GII transposition flow.}
\end{figure}

\end{document}

